\chapter{Multi-Manager Platforms}\label{ch:in:multi-manager-platforms}

In January 2016 the largest multi-manager platform's registered adviser reported 1\,625 employees to the SEC; ten years
later it reported 6\,140, and its home page counted more than 360 investment teams. A second platform grew from 287
employees to 2\,063 over the same decade. No part of the industry has hired as fast. For the people it hires --- portfolio
managers, analysts, researchers, and the engineers and risk managers around them --- a platform is a job with a known
exit condition: a team that loses more than a set fraction of the capital it was given is cut, and then closed. This
chapter describes the platforms as employers, reads their growth in their filings, and measures what a drawdown rule
means for how long a team lasts.

\section{The platform from the employee's side}

Book~16 (chapter~3) described the platform's economics: pods that each run a share of one fund's capital, costs passed
through to investors, payouts on each pod's own profit, a centre book that hedges what the pods hold in common. From
the employee's side the same structure looks like this.
\begin{itemize}
\item \textbf{A team is a small business inside a large one.} A portfolio manager hires analysts and a researcher or
two, agrees a capital allocation and risk limits, and is paid a share of the team's own profit after its costs.
\item \textbf{Everything else is provided.} Data, execution, technology, risk systems, compliance, operations and
financing are the platform's; a team uses them and pays for them through its costs.
\item \textbf{The limits are the contract.} A team's drawdown from its peak, measured on the capital it was allocated,
triggers a cut in capital and then the team's closure (the de-risking ladder of Book~16, chapter~8).
\item \textbf{The platform employs many people who are not in teams}: the risk, technology and data staff who run the
common machinery, and the centre-book traders.
\end{itemize}

The platforms named in this chapter describe themselves in these terms: a firm with a ``multi-strategy, multi-PM
approach'' whose ``165 investment teams span six strategies''; one that has grown ``from a proprietary trading firm with
deep systematic roots into a multi-strategy, multi-manager hedge fund''; ``a global, multi-strategy, multi-manager
investment firm''.

\section{Size: assets, employees and teams}

The platforms' US advisers file Form ADV (chapter~4), and the SEC's files of past months keep the answers, so a platform's
growth can be read from its own filings.

\begin{dated}{2026-01}{Four platforms in their filings and on their pages}\label{dat:in:multi-manager-platforms:size}
\footnotesize
\begin{tabular}{@{}lrrrrr@{}}
\toprule
registered adviser & Jan 2016 & Jan 2017 & Dec 2022 & Dec 2024 & Jan 2026\\
\midrule
Millennium Management, employees & 1\,625 & 1\,830 & 3\,985 & 5\,550 & 6\,140\\
\quad RAUM (\$ billion) & 181.5 & 207.9 & 341.0 & 505.9 & 571.1\\
Balyasny Asset Management, employees & 287 & 395 & 1\,048 & 1\,857 & 2\,063\\
\quad RAUM (\$ billion) & 29.0 & 36.0 & 181.0 & 248.0 & 265.2\\
Schonfeld Strategic Advisors, employees & -- & 89 & 620 & 952 & 872\\
ExodusPoint Capital Management, employees & -- & -- & 688 & 647 & 634\\
\bottomrule
\end{tabular}

\smallskip\noindent The firms' own pages: Millennium, ``7\,000+ employees globally'', ``360+ investment teams'', ``\$97BN+
AUM''; Balyasny, ``165 investment teams'' in six strategies, ``2000+ investment team and support professionals'',
``\$37B assets under management''. RAUM is gross (chapter~4); the firms' AUM figures are net.
\end{dated}

The filings tell three things the firms' pages do not
(\cref{fig:in:multi-manager-platforms:growth}). Growth was fast but not uniform: the largest adviser's staff grew 3.8
times in ten years, a second's 7.2 times; a younger platform's staff fell slightly after 2022. Assets grew about as fast
as staff: the largest adviser's RAUM per employee fell by 17\% over the decade, the second's rose by 27\%. And the share of
staff in advisory functions fell at both, from 63\% to 47\% and from 64\% to 42\%: the platforms built their common
machinery faster than their teams. For a candidate, that is the shape of the job market inside a platform: more of
the new positions are in technology, data and risk than in the teams.

\begin{omfigure}
\begin{tikzpicture}
\begin{axis}[width=11cm,height=5.6cm,xmin=2015.5,xmax=2026.5,ymin=0,ymax=7,xlabel={\small date of the SEC file},
  ylabel={\small employees (thousands)},xtick={2016,2018,2020,2022,2024,2026},xticklabel style={/pgf/number format/1000 sep={}},
  tick label style={font=\footnotesize},grid=major,grid style={gray!20},table/col sep=comma,
  legend style={font=\scriptsize,at={(0.5,-0.3)},anchor=north,legend cell align=left,/tikz/every even column/.append style={column sep=8pt}},legend columns=2]
\addplot[omDef,thick,mark=*] table[x=t,y=employees]{figdata/industry/05-multi-manager-platforms/adv_mlp.csv};
\addplot[omThm,thick,mark=square*] table[x=t,y=employees]{figdata/industry/05-multi-manager-platforms/adv_bam.csv};
\addplot[omProp,thick,mark=triangle*] table[x=t,y=employees]{figdata/industry/05-multi-manager-platforms/adv_sch.csv};
\addplot[gray,thick,mark=diamond*] table[x=t,y=employees]{figdata/industry/05-multi-manager-platforms/adv_exo.csv};
\legend{Millennium Management,Balyasny Asset Management,Schonfeld Strategic Advisors,ExodusPoint Capital Management}
\end{axis}
\end{tikzpicture}

\omcaption{Employees reported on Form ADV by four platforms' US advisers, in the SEC's files of January 2016, January
2017, December 2022, December 2024 and January 2026 (no file between 2017 and 2022 was read; the lines join the points
read). Data: SEC adviser files, item 5.A, through \texttt{in\_platforms.panel}.}\label{fig:in:multi-manager-platforms:growth}
\end{omfigure}

\section{Team turnover: what the public record shows}

How many teams a platform hires and closes each year is the number a portfolio manager most wants and the one no
platform publishes. The press reports hirings and departures of individual teams, and occasionally a total for a year;
none of it is a primary source, and this book does not print it. What can be said from public records is structural.

\begin{definition}[Team tenure, team turnover rate]\label{def:in:multi-manager-platforms:tenure}
A team's \emph{tenure}\index{team tenure} is the time from its start on a platform to its closure or departure. A
platform's \emph{team turnover rate}\index{team turnover rate} is the number of teams that close or leave in a year
divided by the average number of teams running during it.
\end{definition}

The two are linked: if teams left at a constant rate $\lambda$ a year, \omterm{def:in:multi-manager-platforms:tenure}{tenure} would be exponential with median
$\ln 2/\lambda$, and a platform with a turnover rate of 20\% would see half of each year's new teams leave within about
3.5 years. The rate is not constant in practice: a new team is most at risk in its first year, before its gains give it
room below the peak. The next section computes the curve.

\section{Drawdown limits and how long a team lasts}

\begin{method}[Tenure under a drawdown ladder]\label{met:in:multi-manager-platforms:tenure}
\begin{enumerate}
\item Give each team a true Sharpe ratio and a volatility of P\&L on its allocated capital.
\item Run its daily P\&L; halve its capital when its drawdown from the peak since it started exceeds the cut level, and
close it when the drawdown exceeds the stop level.
\item Record the time of closure, or the horizon if it survives.
\item Estimate the survival curve and the median \omterm{def:in:multi-manager-platforms:tenure}{tenure}; for a mix of teams, the turnover rate and the share of closed
teams that were skilled.
\end{enumerate}
\end{method}

The chapter's teams (illustrative, not any platform's terms) run 10\% annual volatility on their capital. Two ladders
are compared: a tight one that halves capital after a 5\% drawdown and closes the team after 7.5\%, and a loose one at
10\% and 15\%. The drawdown is measured since the team started, so a team that has made money has room.

\begin{center}\footnotesize
\begin{tabular}{@{}lrrrrr@{}}
\toprule
ladder, true Sharpe & median \omterm{def:in:multi-manager-platforms:tenure}{tenure} & closed in 10 years & alive at 1 year & at 3 years & turnover\\
\midrule
tight, 0.0 & 0.9 years & 99.6\% & 45.6\% & 15.9\% & 63\%\\
tight, 0.5 & 1.5 years & 94.4\% & 60.2\% & 33.9\% & 33\%\\
tight, 1.0 & 3.6 years & 74.8\% & 72.2\% & 53.5\% & 16\%\\
tight, 1.5 & over 10 years & 47.9\% & 81.8\% & 69.5\% & 7\%\\
loose, 0.5 & over 10 years & 41.4\% & 96.4\% & 79.1\% & 6\%\\
loose, 1.0 & over 10 years & 15.3\% & 98.7\% & 92.8\% & 2\%\\
\bottomrule
\end{tabular}
\end{center}

Under the tight ladder even a good team is more likely than not to be closed within ten years: a team with a Sharpe
ratio of 1.0 has a median \omterm{def:in:multi-manager-platforms:tenure}{tenure} of 3.6 years, and one of 0.5 of 1.5 years
(\cref{fig:in:multi-manager-platforms:surv}). The loose ladder keeps most good teams for a decade. In a platform whose
teams are spread evenly across Sharpe ratios of 0, 0.5, 1.0 and 1.5, the tight ladder closes 20\% of teams a year and
the loose one 5\%; 35\% of teams are closed in their first year under the tight ladder, 3\% under the loose one.

\begin{omfigure}
\begin{tikzpicture}
\begin{axis}[width=11cm,height=5.8cm,xmin=0,xmax=10,ymin=0,ymax=1.02,xlabel={\small years since the team started},
  ylabel={\small share of teams still running},tick label style={font=\footnotesize},grid=major,grid style={gray!20},
  table/col sep=comma,legend style={legend cell align=left,font=\scriptsize,at={(0.5,-0.3)},anchor=north,/tikz/every even column/.append style={column sep=8pt}},legend columns=2]
\addplot[omThm,thick] table[x=years,y=tight05]{figdata/industry/05-multi-manager-platforms/survival.csv};
\addplot[omThm,thick,dashed] table[x=years,y=tight10]{figdata/industry/05-multi-manager-platforms/survival.csv};
\addplot[omDef,thick] table[x=years,y=loose05]{figdata/industry/05-multi-manager-platforms/survival.csv};
\addplot[omDef,thick,dashed] table[x=years,y=loose10]{figdata/industry/05-multi-manager-platforms/survival.csv};
\legend{{tight ladder, Sharpe ratio 0.5},{tight ladder, Sharpe ratio 1.0},{loose ladder, Sharpe ratio 0.5},{loose ladder, Sharpe ratio 1.0}}
\end{axis}
\end{tikzpicture}

\omcaption{Survival of simulated teams under two drawdown ladders: the share still running after each year, 4\,000 teams
per curve, annual volatility 10\% of allocated capital. Tight ladder: capital halved at a 5\% drawdown from the peak since
the start, team closed at 7.5\%; loose ladder: 10\% and 15\%. Illustrative parameters. Data:
\texttt{firm.teamtenure.stop\_times}, through \texttt{in\_platforms.times}.}\label{fig:in:multi-manager-platforms:surv}
\end{omfigure}

The ladder cannot tell skill from luck quickly. Of the teams the tight ladder closes, 39\% had a true Sharpe ratio of
1.0 or more; under the loose ladder, 14\%. The halving step matters: without it, a team with a Sharpe ratio of 0.5 lasts a
median 0.7 years under the tight stop instead of 1.5, because the cut halves the size of the losses that follow a bad
start. A platform chooses where to stand between closing bad teams fast and keeping good ones, and a portfolio manager
joining it should ask where that is.

\begin{remark}[What the model leaves out]\label{rem:in:multi-manager-platforms:model}
Real platforms review teams on more than drawdowns: returns against the allocated risk, crowding with other teams,
behaviour under limits. Teams also leave on their own, for another platform or to start a fund (Book~16, chapter~25).
Capital is reallocated upward after gains. The model isolates the one rule a team signs up to; it is not a forecast of
any platform's turnover, and the teams' volatility and Sharpe ratios are chosen for illustration.
\end{remark}

\section{How platforms hire}

A platform hires a portfolio manager for a record and a plan. The record is the manager's past P\&L with the risk taken
to earn it, which Book~16 (chapter~25) calls portable only if it can be attributed and verified; the plan is the capital
the manager asks for, the strategy, the team and the budget. Analysts and researchers are hired by the portfolio
managers, often from other platforms' teams or from banks' desks; engineers and risk staff by the platform itself.
Because every team's economics are its own, a platform's hiring runs as a market: teams compete for capital and people,
and a team that is closed releases both.

For an analyst the consequence is that the employer is, in practice, the team: when the team closes, the job usually
ends with it, unless another team hires the analyst. \omterm{def:in:multi-manager-platforms:tenure}{Tenure} under a ladder is therefore also the analyst's job security,
and the table above is a reason to ask about it before joining.

\section{Tutorial: growth in the filings and tenure in the model}\label{tut:in:multi-manager-platforms:tut}

\begin{tutorial}
\textbf{Goal.} Read four platforms' growth from their Form ADV filings, and compute \omterm{def:in:multi-manager-platforms:tenure}{team tenure} under two ladders.
\textbf{End state:} the dated box, \cref{fig:in:multi-manager-platforms:growth}, the \omterm{def:in:multi-manager-platforms:tenure}{tenure} table and
\cref{fig:in:multi-manager-platforms:surv}.
\begin{tutsteps}
\item \textbf{The panel.} \texttt{in\_platform\_derive.py} reads five of the SEC's adviser files (2016 to 2026, xlsx and
CSV formats) with \texttt{firm.formadv.read(path, hedge\_only=False)} and keeps the four advisers by their CRD numbers;
\texttt{in\_platforms.growth} compares two dates.
\item \textbf{One team.} \texttt{firm.teamtenure.stop\_times} runs a batch of teams under a ladder
(\cref{lst:in:multi-manager-platforms:stop}); returns come from Book~16's \texttt{firm.podshop.pods}.
\omcode{code/firm/teamtenure/firm_teamtenure.py}{35}{48}{The ladder applied day by day: halve at the cut, close at the
stop, both measured from the peak since the team started.}\label{lst:in:multi-manager-platforms:stop}
\item \textbf{The table.} \texttt{in\_platforms.table()} gives the median \omterm{def:in:multi-manager-platforms:tenure}{tenure}, the share closed, survival at one and
three years and the turnover rate for each ladder and Sharpe ratio; \texttt{mixed()} the platform-level figures.
\item \textbf{Check the model.} Run the tight ladder with four times as many steps a year: the median \omterm{def:in:multi-manager-platforms:tenure}{tenure} of a team
with a Sharpe ratio of 0.5 stays within 10\% of 1.5 years. Remove the halving step: it falls to 0.7 years.
\end{tutsteps}
\end{tutorial}

\textbf{What to change next.} Raise the teams' volatility to 15\% and see how the tight ladder's turnover changes
(exercise~7); give teams a Sharpe ratio that falls to zero after two years and ask how fast each ladder closes them.

\section{Build: the tenure model}\label{bld:in:multi-manager-platforms:teamtenure}

\begin{build}
\bfield{Purpose} Measure what a drawdown rule does to how long teams last, for this chapter and for chapter~22's
portfolio-manager deal.
\bfield{Interface} \texttt{firm.teamtenure}: \texttt{Ladder(cut, stop)}; \texttt{stop\_times(sr, vol, ladder, n, years,
rng, days)}; \texttt{survival(times, grid)}; \texttt{median\_tenure(times)}; \texttt{turnover(times, years)};
\texttt{false\_cut\_share(times\_by\_sr, skilled)}. Wraps \texttt{firm.podshop.pods}.
\bfield{Rules} Drawdowns are measured on the capital actually run, from the peak since the start; a closed team stays
closed; teams alive at the horizon are censored, not closed.
\bfield{Acceptance tests} \texttt{code/firm/teamtenure/tests/}: survival, median and turnover on known times; more skill
and a looser ladder mean fewer closures; the false-cut share on constructed times.
\bfield{Stretch} Capital reallocation after gains; teams that leave on their own at a constant rate; a Sharpe ratio that
decays with the team's age.
\end{build}

\begin{omsources}
\item SEC, Information about registered investment advisers: monthly files of January 2016, January 2017, December 2022,
December 2024 and January 2026 (Internet Archive copies).
\item The platforms' own pages: Millennium, Balyasny, Schonfeld, ExodusPoint (ledger F2--F5).
\item Book~16, chapters 3, 8 and 25, for the platform's economics, the de-risking ladder and track records.
\end{omsources}

\section{Exercises}

\begin{exercise}[$\star$]\label{exo:in:multi-manager-platforms:1}
From the dated box, compute the growth of the largest adviser's employees and RAUM between January 2016 and January 2026.
\end{exercise}

\begin{exercise}[$\star$]\label{exo:in:multi-manager-platforms:2}
If teams leave a platform at a constant 20\% a year, what is the median \omterm{def:in:multi-manager-platforms:tenure}{tenure} of a team?
\end{exercise}

\begin{exercise}[$\star$]\label{exo:in:multi-manager-platforms:3}
A platform's page counts 360 teams and 7\,000 employees. What does that suggest about the share of staff outside the
teams, if a team averages five people?
\end{exercise}

\begin{exercise}[$\star\star$]\label{exo:in:multi-manager-platforms:4}
Compute the advisory share of the second platform's adviser in January 2016 and January 2026, and say what the change
means for a candidate.
\end{exercise}

\begin{exercise}[$\star\star$]\label{exo:in:multi-manager-platforms:5}
Under the tight ladder, what share of teams with a Sharpe ratio of 1.0 are still running after three years? After how
long is half of them closed?
\end{exercise}

\begin{exercise}[$\star\star$]\label{exo:in:multi-manager-platforms:6}
Read \cref{fig:in:multi-manager-platforms:surv}: at which age does the gap between the tight and loose ladders for teams
with a Sharpe ratio of 0.5 exceed 40 percentage points?
\end{exercise}

\begin{exercise}[$\star\star\star$]\label{exo:in:multi-manager-platforms:7}
\emph{Coding.} Rerun \texttt{stop\_times} for a Sharpe ratio of 1.0 under the tight ladder with a volatility of 15\%
instead of 10\%. What happens to the median \omterm{def:in:multi-manager-platforms:tenure}{tenure}, and why?
\end{exercise}

\begin{exercise}[$\star\star\star$]\label{exo:in:multi-manager-platforms:8}
\emph{Find the flaw.} ``Under the tight ladder 39\% of closed teams were skilled, so the platform should stop closing
teams.''
\end{exercise}

\section{Problem: How Long Does a Team Last?}

\begin{problem}[{Weekend problem --- how long does a team last?}]\label{pb:in:multi-manager-platforms:1}
A portfolio manager is offered a seat at a platform with a drawdown ladder. She believes her strategy's Sharpe ratio is
between 0.5 and 1.0 and wants to know her chances.

\medskip\noindent\textbf{Part I --- The platform.}
\begin{enumerate}
\item List what a team provides and what the platform provides.
\item Give the largest adviser's employees in January 2016 and January 2026, and the growth factor.
\item Give the change in its RAUM per employee and in its advisory share over the decade.
\item What do the two platforms' pages say about their number of teams?
\item Why does the chapter not print the press's counts of team closures?
\end{enumerate}

\medskip\noindent\textbf{Part II --- \omterm{def:in:multi-manager-platforms:tenure}{Tenure}.}
\begin{enumerate}[resume]
\item Define \omterm{def:in:multi-manager-platforms:tenure}{team tenure} and the turnover rate, and relate the median \omterm{def:in:multi-manager-platforms:tenure}{tenure} to a constant rate $\lambda$.
\item State the two ladders and the teams' volatility.
\item Give the median \omterm{def:in:multi-manager-platforms:tenure}{tenure} for Sharpe ratios 0.5 and 1.0 under the tight ladder.
\item Give the share of such teams closed within ten years under each ladder.
\item Why is a new team most at risk in its first year?
\end{enumerate}

\medskip\noindent\textbf{Part III --- The platform's view.}
\begin{enumerate}[resume]
\item Give the turnover rate of the mixed platform under each ladder.
\item Give the share of teams closed in their first year under each ladder.
\item Give the share of closed teams that were skilled (Sharpe ratio 1.0 or more) under each ladder.
\item What does the halving step do to the \omterm{def:in:multi-manager-platforms:tenure}{tenure} of a team with a Sharpe ratio of 0.5?
\item Check the model: what happens to that median when the step is divided by four?
\end{enumerate}

\medskip\noindent\textbf{Part IV --- The verdict.}
\begin{enumerate}[resume]
\item State the \emph{named result}: median \omterm{def:in:multi-manager-platforms:tenure}{tenure} under the tight ladder for Sharpe ratios of 0.5 and 1.0, and the
share of closed teams that were skilled under each ladder.
\item What three things should she ask the platform before joining?
\item Why is the \omterm{def:in:multi-manager-platforms:tenure}{tenure} also her analysts' job security?
\item What does the model leave out that could lengthen or shorten her \omterm{def:in:multi-manager-platforms:tenure}{tenure}?
\item In two sentences, what should she conclude?
\end{enumerate}
\end{problem}

\section{Interview questions}

\begin{interviewq}[$\star$ \iqroles{trader, researcher}]\label{iq:in:multi-manager-platforms:1}
What is a multi-manager platform, and how is working in one of its teams different from working at a single-strategy
fund?
\end{interviewq}

\begin{interviewq}[$\star$ \iqroles{researcher, risk}]\label{iq:in:multi-manager-platforms:2}
A team runs 10\% annual volatility with a Sharpe ratio of 1. What is the probability that it is down 5\% or more after
its first year?
\end{interviewq}

\begin{interviewq}[$\star\star$ \iqroles{risk}]\label{iq:in:multi-manager-platforms:3}
Why does a platform halve a team's capital after a first drawdown rather than close it at once?
\end{interviewq}

\begin{interviewq}[$\star\star$ \iqroles{researcher}]\label{iq:in:multi-manager-platforms:4}
A platform closes 20\% of its teams every year. What can you infer about the skill of the teams that remain?
\end{interviewq}

\begin{interviewq}[$\star\star$ \iqroles{trader}]\label{iq:in:multi-manager-platforms:5}
As a new portfolio manager, how would you size risk in your first months under a drawdown ladder, and why?
\end{interviewq}

\begin{interviewq}[$\star\star\star$ \iqroles{researcher, risk}]\label{iq:in:multi-manager-platforms:6}
A team's P\&L is a Brownian motion with drift $\mu>0$ and volatility $\sigma$. Show that a stop at a loss $d$ below
its starting capital catches it with probability $e^{-2\mu d/\sigma^2}$, while a stop at $d$ below its running peak catches
it with probability one. What does that mean for a ladder measured from the peak?
\end{interviewq}
